\exercise{Jacobian derivation\label{exStripesJac}}{4}
In the chapter we used a system-level approach to compute the Jacobian of a general system where the dynamics of variable $n$ on node $i$ is 
\eqn{
\dot{x}_{ni}=f_n(x_{:i})-c_n\sum L_{ij} x_{nj}
}

\subquestion 
Compute Jacobian elements using the node-level approach
\eqn{
{\bf J}_{ni,mj} = \left. \frac{\partial \dot{x}_ni}{\partial x_{mj}} \right|_*
}
and show that it can still be written as ${\bf J}={\bf I}\otimes {\bf P}-{\bf L}\otimes{\bf C}$

\solution
We first differentiate the first term of $\dot{x}_{ni}$, 
\eq{
 \frac{\partial}{\partial x_{mj}} f_n(x_{:i}) = \delta_{ij} \frac{\partial}{\partial x_{mi}} f_n(x_{:i}) = \delta_{ij}P_{nm}   
}
In the first step we noticed that the derives are all with respect to variables in node $j$ while the $f(x_{:i})$ only depends on variables in node $i$ hence the result is zero unless $i=j$, which we express using the Kronecker delta. With the $\delta_{ij}$ factor in place we can then replace the $j$ by a $i$. We are allowed to do this because either $i=j$ or the result is 0 due to the delta anyway. After the exchange we are left with the derivative of function $f_n$ with respect to the variable $x_m$ in the same node, which we recognize as the element $P_{nm}$ of the node Jacobian.  

Now for the second term of $\dot{x}_{ni}$, where we need to rename the summation index to $k$ to avoid confusion 
\eqa{
 \frac{\partial}{\partial x_{mj}} c_n\sum_k L_{ik} x_{nk} 
 &=& c_n\sum_k L_{ik} \frac{\partial}{\partial x_{mj}}  x_{nk} \\ 
 &=& c_n\sum_k L_{ik} \delta_{nm} \delta{jk} \\  
 &=& c_n\delta_{nm} L_{ij} \\  
}
Hence our derivative is 
\eq{
{\bf J}_{ni,mj} =  \delta_{ij}P_{nm} - c_n\delta_{nm} L_{ij}
}
To write this in matrix form we recall that $i,j$ denote the block corresponding to one node while $n,m$ are the indices in the block. 
So let's write an equation for the Jacobian block $i,j$. This means we want to replace the terms with indices $n,m$ with the corresponding matrices. The first term contributes the element $P_{nm}$ to element $n,m$ of our block, so this is just the matrix $\bf P$. The second term gives us the element $c_n \delta_{nm}$, so this is a matrix that has only elements on the diagonal due to the $\delta_{nm}$ and to the diagonal element $n,n$ it contributes $c_n$. So this is our coupling matrix 
\eq{
{\bf C} = {\rm diag}(c_1,c_2, \ldots) = \aveccc{c_1 & & \\ & c_2 & \\ & & \ddots}
}
Using $\bf P$ and $\bf C$ we can now write 
\eq{
{\bf J_{ij}} = \delta_{ij} {\bf P} - L_{ij} {\bf C}
}
This gives us an equation for the block $i,j$ of the Jacobian but we want an equation for the entire Jacobian. We recognize products such as $\delta_{ij} {\bf P}$ now as a block-wise notation of the Kronecker product. For example we could write
\eq{
({\bf A} \otimes {\bf B} )_{ij} = A_{ij} {\bf B} 
}
Now we do this step backwards. The first term contains a $\delta_{ij}$ which describes a matrix that has ones for $i=j$ and zeros elsewhere, i.e. the identity matrix $\bf I$, wheras the second term contains the elements of the Laplacian, so stepping from the blockwise to the system level we can write 
\eq{
{\bf J} = {\bf I} \otimes {\bf P} - {\bf L} \otimes {\bf C} 
}
which is the desired result. 

\subquestion 
Compute the Jacobian again using the intermediate blockwise approach
\eqn{
{\bf J_{ij}} = \left. \frac{\partial \vec{\dot{x}_{:i}}}{\partial \vec{x_{:j}}} \right|_*.
}

\solution 
We define 
\eq{
\vec{x_{:i}}=\avec{x_{1,i}\\x_{2,i}\\ \vdots \\ x_{N,i}}
}
The dynamics of this vector can be written as 
\eq{
\vec{\dot{x}_{:i}} = f(\vec{x_{:i}}) - {\bf C} \sum_k L_{ik}  \vec{x_{:k}}
}
where we again identified $\bf C$ and changed the summation index to $k$ as above. Importantly all indices $n$ and $m$ have now vanished since we are writing equations for block vectors and hence don't need to resolve within-block coordinates anymore. Note in particular, that the equation now contains the function $f$ which returns the vector of changes due to the local dynamics within a nodes, in contrast to $f_n$ which returns only the change in the $n$-th variable of the node. (In case you are wondering this is not a new notation, we have been using it ever since we write the equation $\dot{\vec{x}}=f(\vec{x})$ ten chapters ago). 

We can now compute the Jacobian blocks as 
\eqa{
{\bf J_{ij}} &=& \left. \frac{\partial \vec{\dot{x}_{:i}}}{\partial \vec{x_{:j}}} \right|_*  \\
 &=& \left. \frac{\partial}{\partial \vec{x_{:j}}} f(\vec{x_{:i}}) - {\bf C} \sum_k L_{ik}  \vec{x_{:k}} \right|_*  \\
  &=& \left. \frac{\partial}{\partial \vec{x_{:j}}} f(\vec{x_{:i}}) \right|_* -  \left. {\bf C} \sum_k L_{ik} \frac{\partial}{\partial \vec{x_{:j}}} \vec{x_{:k}} \right|_*  \\
  &=& \left. \delta_{ij} \frac{\partial}{\partial \vec{x_{:i}}} f(\vec{x_{:i}}) \right|_* -  \left. {\bf C} \sum_k L_{ik} \delta_{jk}   \right|_*  \\
  &=& \delta_{ij} {\bf P}  - L_{ij} {\bf C}    
}
We have now arrived at an equation that we recognize from part a) of this exercise. Repeating the same step as before leads to
\eq{
{\bf J} = {\bf I} \otimes  {\bf P} - {\bf L} \otimes {\bf C}
}
These sort of calculations are quite important as they allow us to recongize the deeper structure of complex equation systems (and it won't always be the same structure as in this chapter). These sort of calculations are not particularly easy, but with some practice they become manageable. Generally the micro-level calculation with the multi-index notation offers the most easily navigable, but also often perhaps most tedious approach. The meso-scale block-level derivation, can be faster but may require more navigational skills. Finally the system level derivation, where we compute the entire Jacobian in one go can be extremely elegant, but also extremely difficult to navigate.  
